\clearpage\thispagestyle{empty}
\pdfbookmark[0]{Sources for Extracts from Gauss’s September 1844 Comet Letters}{cometsep1844-sources}
\begin{center}\textbf{Sources for Extracts from Gauss’s September 1844 Comet Letters}\end{center}
\noindent This AI translation follows the extracts of 1, 3, 6 and 8 September 1844 on pp.~468--469 of \href{https://gdz.sub.uni-goettingen.de/id/PPN236007467}{\textit{Werke}, vol.~VI} (Göttingen, 1874). Carl Friedrich Gauss wrote to Heinrich Christian Schumacher. The journal printed the four extracts together in \textit{Astronomische Nachrichten}, vol.~XXII, no.~516, \href{https://api.digitale-sammlungen.de/iiif/image/v2/bsb10049447_00103/full/1800,/0/default.jpg}{cols.~189--190} and \href{https://api.digitale-sammlungen.de/iiif/image/v2/bsb10049447_00104/full/1800,/0/default.jpg}{cols.~191--192}. The collected citation dates the issue to 3 October 1844.
\par\smallskip\noindent The collected 31 August right ascension has a raised \(u\) after the 8, where the journal uses a double prime. The collected 6 September extract describes an observed result that takes right ascension into account. The first printing describes a calculated result that takes refraction into account. The translation follows the collected readings.
\par\smallskip\noindent Gauss’s remark on the radius vector refers to the coordinate formulas in \href{https://api.digitale-sammlungen.de/iiif/image/v2/bsb10049447_00066/full/1800,/0/default.jpg}{no.~512, col.~116} of the journal, cited in the \hyperlink{native-472-notice}{collected notice of Goldschmidt’s elements}. Translation and typesetting produced by AI.
\clearpage
